\NoteChapter{5}{Chapter 6 紧算子与谱分解定理}
\NoteSection{1}{基本知识}


\begin{definition}[紧算子]
    \label{def:compact_operator}
    设$E,F$是Banach空间.
    有界线性算子$T : \mathcal{L}(E,F)$若满足$T(B_E)$是$F$中的相对紧集（闭包为$F$中紧集），
    则$T$ 
    称为紧算子.
    $E \to F$的所有紧算子的集合记作$\mathcal{K}(E,F)$.
\end{definition}

\begin{remarkbox}
\textbf{Remark}：紧算子定义的本质是，将有界集映为相对紧集.
设$A$是$E$中的有界集，则存在$M>0$，使得$A \subset MB_E$，
因此$T(A) \subset MT(B_E)$.
相对紧集的子集仍然是相对紧集，
因此$T(A)$是$F$中的相对紧集.
\end{remarkbox}

\begin{theorem}
    $\mathcal{K}(E,F)$是$\mathcal{L}(E,F)$的闭子空间（按算子范数）.
\end{theorem}


\begin{proof}
    子空间是显然的.
    设$T_n \in \mathcal{K}(E,F)$，且$T_n \to T \in \mathcal{L}(E,F)$.
    我们要证明$T$是紧算子.
    即证$T(B_E)$是$F$中的相对紧集.

    已知$T_n(B_E)$是相对紧集，
    因此是完全有界集.
    对于任意$\varepsilon>0$，存在有限个点$\{y_1,\dots,y_m\} \subset T_n(B_E)$，使得
    \[
        T_n(B_E) \subset \bigcup_{j=1}^m B_F(y_j,\varepsilon).
    \]
    由于$T_n \to T$，存在$N>0$，当$n>N$时，有$\|T_n-T\|<\varepsilon$.
    也就是说，
    对于任意$x \in B_E$,
    $|T(x) - T_n(x)| \le \varepsilon \|x\| \le \varepsilon$.
    因此
    \[
        T(B_E) \subset T_n(B_E) + B_F(0,\varepsilon)
        \subset \bigcup_{j=1}^m B_F(y_j,2\varepsilon),
    \]
    这说明了$T(B_E)$完全有界集，因此是相对紧集.
    因此$T\in\mathcal{K}(E,F)$. 
\end{proof}

    \begin{definition}[有限秩算子]
        \label{def:finite_rank_operator}
        设$E,F$是Banach空间.
        若算子$T : E \to F$的值域$T(E)$是有限维的，
        则称$T$为有限秩算子.
    \end{definition}

    \begin{remarkbox}
    \textbf{Remark}：有限秩算子是有限维空间线性变换的最直接推广，
    而紧算子则是更重要的推广.
    有限秩算子一定是紧算子.
    这是因为，
    若$T$是有限秩算子，
    则$T(B_E)$是有限维空间中的有界集，
    因此是相对紧集.
    \end{remarkbox}

    \begin{corollary}
        设$\{T_n\}$是一列有限秩算子，
        $T_n \to T$，则$T$是紧算子.
    \end{corollary}

    \begin{theorem}[Schauder]
        若$T \in \mathcal{K}(E,F)$，
        则其共轭算子$T^\star \in \mathcal{K}(F^\star,E^\star)$.
    \end{theorem}

    \begin{proof}
            设$B_{F^\star}$是$F^\star$中的单位闭球.
    我们要证明$T^\star(B_{F^\star})$是$E^\star$中的相对紧集.
    任取$B_{F^\star}$中点列$\{v_n\}$，
    去证$T^\star(v_n)$存在收敛子列.

    记$K = \overline{T(B_E)}$是紧集.
    将$\{v_n\}$看成其在$K$上的限制，
    它是$C(K)$中的点列，
    而且一致有界.
    由Arzelà–Ascoli定理，
    存在子列$\{v_{n_k}\}$一致收敛到某个连续映射$v \in C(K)$，即

    \[
        \sup_{y \in K} |v_{n_k}(y) - v(y)| \to 0.
    \]

    将$y$替换为$T(u), u \in B_E$

    \[
        \sup_{u \in B_E} |v_{n_k}(T(u)) - v(T(u))| \to 0.
    \]

    根据共轭算子的定义，可改写成

    \[
        \sup_{u \in B_E} |T^\star(v_{n_k})(u) - T^\star(v)(u)| \to 0.
    \]

    这说明$T^\star(v_{n_k})$在$E^\star$中（依算子范数）收敛到$T^\star(v)$.
    因此$T^\star(B_{F^\star})$是相对紧集.   
    \end{proof}


\NoteSection{2}{Riesz-Fredholm理论}

我们有以下两个熟知的结论：

\begin{lemma}[Riesz引理]
    设$E$是赋范线性空间，$M$是$E$的闭子空间，且$M \neq E$.
    则对于任意$\varepsilon \in (0,1)$，存在$x_\varepsilon \in E$，使得
    \[
        \|x_\varepsilon\| = 1,
        \quad
        d(x_\varepsilon,M) \ge 1 - \varepsilon.
    \]
\end{lemma}

\begin{theorem}[Riesz]
    设$E$是赋范线性空间，$B_E$是紧集，
    则$E$是有限维的.
\end{theorem}

在有限维线性空间中，
设$A$是线性算子，
线性方程$Ax = b$的解与$A^\star$的性质密切相关：

(1) 若$Ax= 0$只有零解，则$A^\star x = 0$也只有零解.

(2) 若$Ax = 0$有非零解，则$A^\star x = 0$也有非零解.
方程组$Ax = b$有解当且仅当$b$与$A^\star x = 0$的解正交.

我们希望把这些结论推广到无限维空间中的紧线性算子$T$上：


\begin{theorem}[Fredholm择一定理]
    设$T \in \mathcal{K}(E)$，则：

    (a) $N(I-T)$是有限维的；

    (b) $R(I-T)$是闭的；$R(I-T) = N((I-T)^\star)^{\perp}$；

    (c) $N(I-T) = \{0\}$当且仅当$R(I-T) = E$；

    (d) $\dim N(I-T) = \dim N((I-T)^\star)$.
\end{theorem}

\begin{proof}
    
    (a) 记$E_1 = N(I-T)$，
    对于$u \in B_{E_1} \subset B_E$，
    有$u = Tu \in T(B_E)$.
    因此，$B_{E_1} \subset T(B_E)$.
    由于$T$是紧算子，$T(B_E)$是相对紧集，
    因此$B_{E_1}$是相对紧集.
    由Riesz定理，$E_1$是有限维的.

    \medskip

    (b) 设$f_n = u_n - T(u_n) \to f$，
    要证$f \in R(I-T)$.
    即：存在$g \in E$，使得$g - T(g) = f$.

    记$d_n = d(u_n , N(I-T))$.
    因此$N(I-T)$有限维，
    存在$v_n \in N(I-T)$使得$d_n = \|u_n - v_n\|$.

    先用反证法证明$\{d_n\}$有界（注意$\|u_n\|$本身不一定有界）.
    假设它无界，
    则存在子列$d_{n_k}$，
    使得$d_{n_k} \to \infty$.
    这时候考虑$u_n - v_n$对应的单位向量$w_n = \dfrac{u_n - v_n}{\|u_n - v_n\|}$的收敛性.
    一方面，因为
    \[
    f_n = (u_n - v_n) - T(u_n - v_n) \to f
    \]

    所以$w_n - T(w_n) \to 0$.
    也就是说$d (w_n , N(I-T)) \to 0$.
    但另一方面，按$w_n$的定义，$d (w_n , N(I-T)) = 1$. 矛盾.

    现在知道了$\{d_n\}$有界.
    由于$T$是紧算子，$T(u_n - v_n)$有收敛子列.
    设$T(u_{n_k} - v_{n_k}) \to \ell$.
    又因为$(u_n - v_n) - T(u_n - v_n) \to f$，
    所以$(u_n - v_n) \to f+\ell$. 于是$f + \ell$就是我们一开始想要的$g$.

    证明了$R(I-T)$是闭子空间，我们就有：
    \[
    R(I-T) = N(I-T^\star)^\perp \quad R(I-T^\star) = N(I - T)^\perp
    \]

    \medskip

    (c) $\Rightarrow$：反证法.
    假设$I-T$是单射，但不是满射，
    设$E_n = (I-T)^n (E)$.
    则$\{E_n\}$是严格单调递减的集合列.
    由Riesz引理，
    存在$u_n \in E_n$，
    使得$\|u_n\| = 1$且$d(u_n, E_{n+1}) > 1/2$.
    则
    \[
    T (u_n - u_m) = -(I-T) (u_n - u_m) + (u_n - u_m) 
    \]

    不妨设$n> m$，
    则$-(I-T) (u_n - u_m) + u_n \in E_{m+1}$，
    $\| T (u_n - u_m)\| \ge d (u_m , E_{m+1}) > 1/2$.
    这与$T$是紧算子矛盾.

    $\Leftarrow$：
    若$R(I-T) = E$，
    则$N(I-T^\star) = R(I-T)^\perp = \{0\}$.
    利用$\Rightarrow$的结论，
    $R(I-T^\star) = E^\star$.
    于是
    $N(I-T) = R(I-T^\star)^\perp = \{0\}$.
\end{proof}


\NoteSection{3}{紧算子的谱}

\begin{definition}[预解集]
    设$T \in \mathcal{L}(E)$，
    则$T$的预解集$\rho (T)$定义为
    \[
        \rho(T) = \{\lambda \in \mathbb{R} : \lambda I - T \text{是双射}\}.
    \]

    预解集的补集$\sigma (T) = E \backslash \rho(T)$称为$T$的谱集$\sigma(T)$.
    谱集中的元素若还满足$N(\lambda I - T) = \{0\}$，则称为特征值.
    特征值的集合记作$EV(T)$.

\end{definition}

\begin{proposition}
    设$T$是有界线性算子，
    则$T$的谱集$\sigma(T)$是紧的，
    且$\sigma(T) \subset \{ \lambda \in \mathbb{R} : |\lambda| \leq \|T\| \}$.
\end{proposition}
\begin{proof}
    任取$\lambda > \|T\|$，
    去证$\lambda I - T$是双射.
    任取$f \in E$，
    考虑方程$(\lambda I - T)u = f$.
    整理，
    得到$u = \frac{1}{\lambda}(f + T(u))$.
    注意$u \mapsto \frac{1}{\lambda}(f + T(u))$是压缩映射，
    根据Banach不动点定理，
    方程存在唯一的解$u \in E$.
    这就证明了当$\lambda > \|T\|$时，
    $\lambda I - T$是双射.

    再证$\rho (T)$是开集.
    任取$\lambda_0 \in \rho (T)$.
    对充分小的$\epsilon$，
    以及实数$\lambda$满足$|\lambda - \lambda_0| < \varepsilon$，
    我们希望证明下述方程有唯一解：
    \[
     (\lambda I - T) u = f
    \]

    这等价于
    \[
    (\lambda_0 I - T) u = f + (\lambda_0 - \lambda) u
    \]

    由于$\lambda_0 \in \rho (T)$，
    $\lambda_0 I - T$是双射，
    因此上式等价于
    \[
    u = (\lambda_0 I - T)^{-1} f + (\lambda_0 - \lambda) (\lambda_0 I - T)^{-1} u
    \]

    因为$|\lambda - \lambda_0| = \varepsilon$是充分小的，
    所以$u \mapsto (\lambda_0 I - T)^{-1} f + (\lambda_0 - \lambda) (\lambda_0 I - T)^{-1} u$是压缩映射，
    根据Banach不动点定理，
    方程存在唯一的解$u \in E$.
    这就证明了$\rho (T)$是开集.
\end{proof}

\begin{remarkbox}
    \textbf{Remark}：
    对于一般的线性算子$T$，
    谱点不一定是特征值.
    也就是说，
    有可能$N(\lambda I - T) = \{0\}$，
    但$R(\lambda I - T) \ne E$.
    不过，当$T$是紧线性算子，
    情况就很简单，并且类似于有限维空间线性算子.

    \smallskip

    首先，$0$一定是谱点.
    否则，$T = T - 0I$是双射，
    $I = T \circ T^{-1}$是紧算子，
    是不可能的.
    不过，$0$不一定是特征值点.

    \smallskip

    非零的谱点一定是特征值.
    否则，$N(T - \lambda I) = \{0\}$，
    根据紧算子的性质（Fredholm定理），
    $R(T - \lambda I) = E$，
    但这样$T - \lambda I$就是双射，
    $\lambda \in \rho (T)$，
    矛盾.
\end{remarkbox}

\begin{lemma}
    设$T \in \mathcal{K}(E)$，
    $\{\lambda_n\}$是一列互不相同的实数，$\lambda_n \to \lambda$.
    若$\lambda_n \in \sigma(T) \backslash \{0\}$，
    则$\lambda = 0$.
\end{lemma}

\begin{proof}
    $\lambda_n$是$T$的一些特征值.
    仿照有限维空间的情形，
    可以证明，
    属于不同特征值的特征元素线性无关.
    因此，
    若取每个$\lambda_n$的各一个特征元素$e_n$，
    记$E_n = \mathrm{span} \{ e_1 , e_2, \cdots , e_n \}$.
    则$\{E_n \}$是严格单调递增的有限维子空间列.

    借鉴Fredholm择一定理中(c)的构造方法，
    假设$\lambda \ne 0$，
    我们想要利用$T$的紧性推出矛盾.
    由Riesz引理，
    存在$u_n \in E_n$，
    使得$\|u_n\| = 1$且$d(u_n, E_{n-1}) > 1/2$.
    因此有

    \[
    T(u_n - u_m) = (T - \lambda_n I)(u_n) - (T - \lambda_m I)(u_m) + \lambda_n u_n - \lambda_m u_m
    \]

    不妨设$n > m$.
    注意$(T - \lambda_n I) (u_n) \in E_{n-1}$，
    $(T - \lambda_n I) (u_m) \in E_{m-1} \subset E_{n-1}$.
    因此

    \[
    \| T(u_n - u_m) \| \ge d(\lambda_n u_n , E_{n-1})
    = |\lambda_n| \cdot d(u_n, E_{n-1}) > |\lambda_n| \cdot 1/2
    \]

    由于$\lambda_n \to \lambda \ne 0$，
    当$n$充分大时，
    有$|\lambda_n| > |\lambda|/2$.
    此时

    \[
    \| T(u_n - u_m) \| \ge |\lambda|/2 \cdot 1/2 = |\lambda|/4
    \]

    这与$T$的紧性矛盾.
\end{proof}

\begin{remarkbox}
    \textbf{Remark}：这个结论说明，
    若$T$是紧算子，
    则$\sigma (T) \backslash \{0\}$要么是有限集，
    要么收敛到$0$.
\end{remarkbox}

\NoteSection{4}{紧自伴算子的谱分解定理}

\begin{theorem}[谱分解定理]
    若$T$是自伴
\end{theorem}


